\documentclass[12pt]{article}
\usepackage[utf8]{inputenc}
\usepackage[T1]{fontenc}
\usepackage[margin=1in]{geometry}
\usepackage{setspace}
\usepackage{xcolor}
\usepackage[colorlinks=true, urlcolor=blue!60!black, linkcolor=blue!60!black]{hyperref}
\setlength{\parindent}{0em}
\setlength{\parskip}{0.45em}
\renewcommand{\baselinestretch}{1.0}
\begin{document}
\begin{center}
{\Large \textit{Printz v. United States}, 521 U.S. 898 (1997)\par}
\vspace{0.3em}
{\normalsize Case Brief --- One-Page Edition\par}
\end{center}

\textbf{Memory Jogger.} The Brady Act case: no commandeering of state executive officers (\href{https://www.law.cornell.edu/supct/html/95-1478.ZO.html}{521 U.S.\ 898}).

\textbf{Facts.} Pending a national instant-check system, the Brady Act required local chief law enforcement officers to make a ``reasonable effort'' to check handgun buyers within five business days. Sheriffs Jay Printz and Richard Mack sued.

\textbf{Procedural History.} The district courts struck the mandate; a divided Ninth Circuit reversed.

\textbf{Votes.} Reversed, 5--4 (Scalia for the Court; O'Connor and Thomas concurring; Stevens, Souter, and Breyer dissenting, with Ginsburg).

\textbf{Holding.} Congress may not command state or local officers to administer or enforce a federal regulatory program. The rule is categorical, with no balancing; ministerial reporting duties were reserved.

\textbf{Reasoning.} Early Congresses imposed duties on state judges, not executive officers; \textit{Federalist} No.\ 27 required only non-obstruction; and conscription would undermine dual sovereignty, political accountability, and the unitary executive. \textit{New York v.\ United States} controlled. O'Connor noted the Spending Clause alternative. Stevens relied on the Supremacy Clause and political safeguards; Souter on \textit{Federalist} No.\ 27; Breyer on European federal systems.

\textbf{Examples.} \emph{Same-Side:} forcing local police to honor immigration detainers. \emph{Opposite-Side:} grants conditioned on voluntary checks. \emph{Fence-Sitter:} mandatory missing-children reports.

\textbf{Critique.} Critics note the thin textual basis and ambiguous history; sanctuary jurisdictions now rely on the doctrine.

\textbf{Key Quotation.} ``The Federal Government may neither issue directives requiring the States to address particular problems, nor command the States' officers, or those of their political subdivisions, to administer or enforce a federal regulatory program. It matters not whether policymaking is involved, and no case by case weighing of the burdens or benefits is necessary; such commands are fundamentally incompatible with our constitutional system of dual sovereignty.'' --- Justice Scalia, 521 U.S.\ at 935.

\end{document}
